\scanpage{20}

2) Pour que $L_{\bullet}$ soit \deleted{projectif} \inserted{flasque (projectif)}, il f. et s. qu'il n'y ait que les facteurs $(\Lambda^{\unsure{j|i}}\Phi_{\ast})^{\mu_i}$ \inserted{(resp. \unsure{$=$} $(\Lambda^{0}\Psi_{\bullet})^{\nu_0}=k_{\bullet}^{\nu_0}$)}. Pour que $L_{\bullet}\otimes_k K$ (comme $K$ module $\frac12$ simplicial, où $K$ est le corps des fractions) soit flasque il f. et suffit qu'il n'y ait que des $(\Lambda^i\Phi_{\ast})^{\mu_i}$, $\Phi_{\ast}(i,n)^{f_{i,\ell}}$ \inserted{$(i\geq1)$} à l'exclusion des $(\Lambda^j\Psi_{\ast})^{\nu_j}$ $(j\geq0)$.

\deleted{Pour} 3) Pour déterminer à équiv. près la catégorie additive des k-modules $\frac12$ simpliciaux str. $\partial$ parfaits (resp str. $\partial$ parfaits) il suffit en principe de déterminer les hom mutuels entre les objets indécomposables

\[
\Lambda^i\Phi,\quad \Phi(i,n),\quad\Lambda^j\Psi_{\ast}\qquad(i\geq1,\ j\geq0)
\]

\inserted{un pour chaque $i$, et chaque diviseur $\ell$}

qu'on peut réduire à deux types si on admet aussi le cas $n=1$,

\[
\Phi(i,n)\qquad\Lambda^j\Psi_{\ast}\qquad(i\geq1,\ j\geq0)
\]

On trouve, en repassant aux complexes de chaînes associés

% archive sheet number: 19
